\documentclass[11pt,a4paper]{article}
\usepackage{fullpage}
\usepackage[T1]{fontenc}
\usepackage[french]{babel} \frenchbsetup{og=«,fg=»}
\usepackage{graphicx} % Required for inserting images
\usepackage{amsthm}
\usepackage{amsmath}
\usepackage{amssymb}
\usepackage{stmaryrd}
\usepackage{proof}
\usepackage{booktabs}


\theoremstyle{definition}
\newtheorem{definition}{Définition}[section]
\newtheorem{theorem}{Théoreme}[section]
\newtheorem{lemma}[theorem]{Lemme}



\title{Étude d'une mesure d'espace pour le lambda calcul}
\author{Ewen DUFOUR\\[1ex]
{\small Encadré par}\\ {\small Thibaut BALABONSKI}}
\date{Stage de M1 MPRI\\Mars--Juillet 2025}
\begin{document}
\begin{figure}[!h]
    \centering
    \begin{align*}
    &\infer[sea_v]{((t^\beta~x^\gamma)^\alpha L)_\pi \rightarrow (t^\beta L_{|t})_{L(x) \cdot \pi}}{}
 &\infer[sea_{\neg v}]{((t^\beta~u^\gamma)^\alpha L)_\pi \rightarrow (t^\beta L_{|t})_{(u^\gamma L_{|u}) \cdot \pi}}{}\\\\
&\infer[\beta_{\neg w}]{((\lambda x .t^\beta)^\alpha L_t)_{(u^\gamma L_u )\cdot \pi} \rightarrow (t^\beta[x/u^\gamma L_u]L_t)_\pi}{x \in fv(t) }
&\infer[\beta_w]{((\lambda x .t^\beta)^\alpha L_t)_{(u^\gamma L_u )\cdot \pi} \rightarrow (t^\beta L_t)_\pi}{x \notin fv(t) }\\\\
&\infer[sub]{x^\alpha[x/u^\beta L]_\pi \rightarrow u^\beta L_{\pi}}{}
\end{align*}
    \caption{Reduction rules of the Space Calculus}
    \label{fig:enter-label}
\end{figure}
\end{document}
